% ==================== 第 4 章 研究方案与技术路线 ====================
\clearpage
\section{研究方案与技术路线}

\subsection{总体研究思路}

本课题采用“结构化输入—约束建模—候选求解—独立校验—评分解释—协同发布”
的研究思路。业务规则首先被划分为硬约束和软目标；CP-SAT 负责搜索可行解
并优化指定目标；独立校验器不复用求解器内部状态，而是直接检查输出课表；
应用层再把课表、评分和变更明细呈现给教务人员。该设计把算法正确性、业务
解释和平台交互分离，便于独立测试和迭代。

\subsection{系统总体架构}

图 \ref{fig:architecture-new} 采用黑白分层结构。外部协同平台不直接控制
求解器，而是通过应用服务层交换结构化数据；独立校验器位于求解结果之后，
对外发布前必须通过复核。

\begin{figure}[htbp]
  \centering
  \begin{tikzpicture}[
    node distance=0.75cm,
    layer/.style={draw=black,fill=white,text=black,rounded corners=1pt,
      minimum width=12.4cm,minimum height=1.05cm,align=center,font=\small},
    item/.style={draw=black,fill=white,text=black,minimum height=0.72cm,
      align=center,font=\small,inner xsep=8pt},
    flow/.style={-{Latex[length=2mm]},draw=black,thick}
  ]
    \node[layer] (interaction) {应用交互层：决策驾驶舱 \quad 周课表 \quad 调课中心 \quad 资源管理};
    \node[layer,below=of interaction] (service) {应用服务层：REST API \quad 数据转换 \quad 方案评分 \quad 差异比较};
    \node[layer,below=of service] (decision) {决策核心层：CP-SAT 约束求解器 \quad 多策略目标 \quad 最小影响重排};
    \node[layer,below=of decision] (validation) {质量保障层：输入校验 \quad 独立课表校验 \quad 不可行诊断};
    \node[layer,below=of validation] (data) {业务数据层：教师 \quad 教室 \quad 班级 \quad 课程请求 \quad 时间槽};
    \node[item,left=0.8cm of service] (feishu) {飞书协同端\\（拟集成）};

    \draw[flow] (data) -- (validation);
    \draw[flow] (validation) -- (decision);
    \draw[flow] (decision) -- (service);
    \draw[flow] (service) -- (interaction);
    \draw[flow,<->] (feishu) -- (service);
  \end{tikzpicture}
  \caption{ClassMind 系统总体架构}
  \label{fig:architecture-new}
\end{figure}

\begin{table}[htbp]
  \centering
  \caption{系统模块及职责}
  \label{tab:modules-new}
  \small
  \begin{tabularx}{\textwidth}{L{2.6cm} L{3.1cm} Y}
    \toprule
    \textbf{模块} & \textbf{主要输入输出} & \textbf{职责} \\
    \midrule
    数据模型 & 业务对象、JSON 数据 & 统一字段定义，承载约束所需的结构化信息 \\
    约束求解器 & 问题实例、策略参数；课表 & 构造决策变量与约束，搜索并优化可行方案 \\
    独立校验器 & 问题实例、输出课表；错误清单 & 复核课次完整性、资源互斥、资质、容量和设备条件 \\
    评分与调课服务 & 课表、基线方案；评分和差异 & 计算统一指标，生成最小影响调课结果 \\
    API 与交互界面 & HTTP 请求；页面和 JSON & 连接算法与用户，展示方案、资源和变更明细 \\
    飞书适配层 & 多维表格、消息和审批事件 & 完成字段映射、方案回写和流程触发（第二阶段） \\
    \bottomrule
  \end{tabularx}
\end{table}

\subsection{约束模型}

设课次集合为 $L$，教师集合为 $T$，教室集合为 $R$，时间槽集合为 $S$。
定义布尔决策变量
\[
  x_{l,t,r,s}=\begin{cases}
  1, & \text{课次 $l$ 安排给教师 $t$、教室 $r$、时间槽 $s$；}\\
  0, & \text{否则。}
  \end{cases}
\]

每个课次必须且只能选择一个可行组合：
\[
  \sum_{t\in T}\sum_{r\in R}\sum_{s\in S}x_{l,t,r,s}=1,
  \qquad \forall l\in L.
\]

对任意时间槽，同一教师、教室和班级最多承担一个课次。例如教师互斥约束为
\[
  \sum_{l\in L}\sum_{r\in R}x_{l,t,r,s}\leq 1,
  \qquad \forall t\in T,\ s\in S.
\]
教师资质、教室容量、设备要求及各对象不可用时间在变量生成阶段预筛选；
不满足条件的组合不创建决策变量，从而缩小搜索空间。

软目标采用加权惩罚形式：
\[
  \min Z = w_p P_{\mathrm{pref}} + w_e P_{\mathrm{resource}}
  + w_b P_{\mathrm{balance}} + w_k P_{\mathrm{compact}}
  + w_c P_{\mathrm{change}}.
\]
不同策略只改变权重，不改变硬约束。跨策略比较时使用统一分项指标，而不直接
比较因权重不同而量纲不同的原始目标值。

\subsection{负载均衡与最小影响调课}

教师日负载均衡采用相对目标课次数的绝对偏差作为线性代理，而不是把它表述
为方差的一阶泰勒近似。设教师 $t$ 在日期 $d$ 的课次数为 $c_{t,d}$，目标
课次数为 $\bar c$，则引入非负辅助变量 $u_{t,d}$，并令
\[
  u_{t,d}\geq c_{t,d}-\bar c,\qquad
  u_{t,d}\geq \bar c-c_{t,d}.
\]
最小化 $\sum_{t,d}u_{t,d}$ 即可降低日负载的绝对偏差。

教师请假时，系统先读取原课表作为基线，再把新增不可用时段写入问题实例，
最后以变更课次数量为主要软目标重新求解。

\begin{algorithm}[htbp]
  \caption{最小影响调课流程}
  \label{alg:min-change-new}
  \begin{algorithmic}[1]
    \Require 原问题 $P$、基线课表 $B$、请假教师 $t^*$、时段 $s^*$
    \Ensure 新课表 $B'$、变更集合 $\Delta$
    \State 将 $s^*$ 加入 $t^*$ 的不可用时段，得到新问题 $P'$
    \State 以 $B$ 为基线，为每个发生变化的课次设置变更惩罚
    \State 求解 $P'$，并优先最小化变更课次数量
    \State 使用独立校验器复核新课表 $B'$
    \State 比较 $B$ 与 $B'$，生成变更集合 $\Delta$
    \State \Return $B',\Delta$
  \end{algorithmic}
\end{algorithm}

\subsection{飞书集成方案}

第二阶段拟以飞书多维表格作为业务数据维护入口：教师、教室、班级、课程和
时间槽分别映射为数据表；触发排课后，适配层读取记录并转换为统一 JSON；
求解通过后把课表和评分回写；教务确认方案后再触发审批和消息卡片。飞书
令牌、权限和回调验签统一在适配层处理，求解器不保存平台凭据。

\subsection{技术路线}

研究按“需求与数据字典—约束模型—原型实现—独立校验—实验评价—飞书集成”
推进。每一阶段均产生可验证交付物，模型和实现通过自动化测试连接，避免
界面演示与算法证据脱节。
